% Requires: \usepackage{listings}
\subsection{Separation of Event Gathering and Operational Logging}
\label{subsec:event-gathering-logging}

PolyGraph distinguishes between \emph{execution events} used for measurement and
profiling, and \emph{operational log messages} intended for developers and
operators.  The two concerns previously shared the same event-logging interface:
an \texttt{EVENT\_*} call constructed structured event data and immediately routed
it through the logging subsystem.  This made instrumentation and ordinary logging
conceptually indistinguishable and, more importantly, meant that instrumentation
remained part of the program even when measurements were not required.

The revised design exposes two independent interfaces. Structured execution
instrumentation is provided by the \texttt{events} subsystem, while ordinary
human-readable messages are provided by the \texttt{logging} subsystem.  Both may
internally use spdlog-backed sinks, but they do not share their public API,
configuration, or lifecycle.

\subsubsection{Execution-event gathering}

Execution events describe measurable points or intervals in the execution of an
algorithm. A point event can attach structured fields to a stable event name:

\begin{lstlisting}[language=C++]
#include "events/Events.hpp"

EVENT_INFO("bfs.frontier.expanded",
           {"depth", depth},
           {"frontier_size", frontier.size()},
           {"neighbors", neighbors.size()});
\end{lstlisting}

Timed intervals use an RAII scope. The scope emits a
\texttt{<name>.start} event when it is created and a \texttt{<name>.end} event
when it leaves scope. Additional fields can be attached to the end event:

\begin{lstlisting}[language=C++]
EVENT_SCOPE_INFO(expand_span,
                 "bfs.frontier.expand",
                 {"depth", depth},
                 {"frontier_size", frontier.size()});

const auto rows = store.get_minimal_neighbors(frontier, depth);

EVENT_SCOPE_INFO_FIELD(expand_span, "rows", rows.size());
\end{lstlisting}

At runtime, enabled events are forwarded to an \texttt{EventCollector}. The
collector applies the configured severity and prefix filters, adds common metadata
such as the sequence number, timestamp, elapsed time, service name and thread ID,
and forwards the resulting event to one or more event sinks. In the benchmark
configuration, the principal sink serializes the event as JSON and writes it to
\texttt{logs/polygraph.events.jsonl}. These JSONL records are the input to the
existing event-log reconstruction and timeline tooling.

Event gathering also has a build-time switch:

\begin{lstlisting}[language=bash]
cmake -S . -B build/release \
  -DPOLYGRAPH_ENABLE_EVENT_GATHERING=OFF
\end{lstlisting}

When \texttt{POLYGRAPH\_ENABLE\_EVENT\_GATHERING=OFF}, every
\texttt{EVENT\_*} and \texttt{EVENT\_SCOPE\_*} macro expands to
\texttt{((void)0)}. Consequently, event names, JSON field construction, field
expressions, collector lookups, timers and scope objects are not present at the
instrumented call site. Event implementation source files are also omitted from
the observability target. Thus a performance-sensitive statement such as

\begin{lstlisting}[language=C++]
EVENT_INFO("bfs.layer.completed",
           {"expensive_field", compute_expensive_statistic()});
\end{lstlisting}

does not evaluate \texttt{compute\_expensive\_statistic()} in a build where event
gathering is disabled. This is stronger than a runtime
\texttt{if (events\_enabled)} check because the instrumentation is removed during
preprocessing and cannot contribute runtime work to the measured algorithm.

When event gathering is compiled in, individual events may still be controlled at
runtime through the \texttt{events} section of the observability configuration.
For example:

\begin{lstlisting}
"events": {
  "enabled": true,
  "minimum_level": "debug",
  "periodic_flush_ms": 1000,
  "selection": {
    "include_prefixes": ["bfs.", "neo4j.", "arango."]
  },
  "sinks": [{
    "name": "events-json",
    "type": "rotating_file",
    "format": "json",
    "path": "logs/polygraph.events.jsonl"
  }]
}
\end{lstlisting}

\subsubsection{Operational logging}

Operational logging is independent from execution-event gathering. It is intended
for diagnostic statements that a developer or operator should read, rather than
for benchmark reconstruction. The interface uses format-style messages:

\begin{lstlisting}[language=C++]
#include "logging/Log.hpp"

LOG_INFO("Starting BFS from {} with max depth {}", source_uuid, max_depth);
LOG_DEBUG("Selected backend {} for {} frontier vertices",
          backend_name, frontier.size());
LOG_WARN("Backend {} returned an empty batch", backend_name);
LOG_ERROR("BFS request failed: {}", error.what());
\end{lstlisting}

These messages do not create structured benchmark events and do not appear in the
event timeline. They are routed through the independently configured application
logger. The default configuration writes informational messages to the console and
writes debug-and-higher messages to a rotating file at
\texttt{logs/polygraph.log}:

\begin{lstlisting}
"logging": {
  "enabled": true,
  "minimum_level": "debug",
  "periodic_flush_ms": 1000,
  "sinks": [
    {
      "name": "console",
      "type": "stdout",
      "minimum_level": "info"
    },
    {
      "name": "application-log",
      "type": "rotating_file",
      "minimum_level": "debug",
      "path": "logs/polygraph.log"
    }
  ]
}
\end{lstlisting}

The distinction is visible at call sites. For example, the BFS executable records
one machine-readable lifecycle event and one human-readable log message for the
same operation:

\begin{lstlisting}[language=C++]
LOG_INFO("BFS completed with {} discovered vertices",
         result.distance.size());

EVENT_INFO("bfs.cli.completed",
           {"vertices", result.distance.size()});
\end{lstlisting}

The first statement is operational output. The second statement belongs to the
measurement data model. Disabling event gathering at compile time removes only the
second statement; ordinary application logging remains available.

\subsubsection{Independent lifecycle}

The two facilities are installed and shut down separately:

\begin{lstlisting}[language=C++]
const auto observability =
    polygraph::cli::parse_observability_config(args);

polygraph::logging::install_logger(observability.logging);
polygraph::events::install_collector(observability.events);

// ... run application ...

polygraph::events::shutdown();
polygraph::logging::shutdown();
\end{lstlisting}

This separation allows PolyGraph to be built with no execution instrumentation for
clean performance measurements while retaining normal operational logging, or to
compile event gathering in and select the desired event set at runtime for detailed
profiling runs.
